\documentclass[12pt]{article}

% ---------- PAGE SETUP ----------
\usepackage[a4paper,margin=1in]{geometry}
\usepackage{microtype}
\raggedbottom

% ---------- PACKAGES ----------
\usepackage{graphicx}
\graphicspath{{figures/}}
\usepackage{float}
\usepackage{amsmath}
\usepackage{amssymb}
\usepackage{booktabs}
\usepackage{cite}
\usepackage[font=small,labelfont=bf]{caption}
\usepackage{titlesec}
\usepackage{hyperref}
\hypersetup{hidelinks}

% ---------- SECTION FORMATTING ----------
\titleformat{\subsection}
  {\normalfont\normalsize\bfseries}
  {\thesubsection}{1em}{}
\titleformat{\subsubsection}
  {\normalfont\normalsize\itshape}
  {\thesubsubsection}{1em}{}

% ---------- PARAGRAPH FORMATTING ----------
\setlength{\parskip}{0.5em}
\setlength{\parindent}{0pt}

% ---------- TITLE ----------
\title{A Conceptual Framework for Early Mental Health Risk Detection Using
Conversational Multimodal Elicitation, Consistency Modeling, and
Statistical Baseline Calibration}
\author{Sengo Basar}
\date{February 2026}

\begin{document}

\maketitle

% =====================================================================
\begin{abstract}
Depression and anxiety frequently remain undetected until they reach an
advanced stage, and existing screening approaches are limited in
complementary ways: clinical instruments are costly and depend on
self-report, while single-modality computational models cannot capture
the distributed nature of mental-health signals. A central difficulty is
the discrepancy between what a person explicitly reports and what their
language, voice, face, and behaviour implicitly reveal.

This paper proposes---but does not yet empirically evaluate---a
conversational multimodal framework that treats that discrepancy as a
first-class measurement rather than as noise. The framework specifies
three components. First, a finite state machine (FSM) interview engine
comprising eight clinical modules, grouped into five thematic domains,
that would elicit explicit ordinal self-reports alongside implicit
linguistic, facial, speech, and behavioural features. Second, a dual
elicitation consistency model that projects implicit representations
into the explicit feature space and quantifies their disagreement as a
consistency score $C$. Third, a statistical baseline layer that
positions session-level risk and response latency relative to a
population distribution through z-score normalization.

We argue on design grounds that $C$ should be treated as a reliability
signal rather than as an additive severity term. A discrepancy measure
has no inherent clinical direction, and if $C$ also enters the feature
vector, an additive term double-counts it. The paper therefore specifies
a reliability-gated decision rule and states four falsifiable hypotheses
for future evaluation. Ethical requirements, including an immediate
escalation pathway for disclosed suicidal ideation, are specified in
full. No experimental results are reported; a complete evaluation
protocol is proposed instead.
\end{abstract}

\noindent\textbf{Keywords:} Multimodal AI, Mental Health, Speech Analysis,
Facial Emotion, Risk Detection, Consistency Modeling, Dual Elicitation,
Response Bias, Statistical Baseline, Anomaly Detection, Z-Score
Normalization, Valence--Arousal, Dimensional Emotion Modeling, Conceptual
Framework

% =====================================================================
\section{Introduction}

Mental health conditions are difficult to detect because their
presentation is heterogeneous, nonlinear, and frequently masked.
According to the World Health Organization, approximately one in eight
people---roughly 970 million individuals in 2019---lived with a mental
disorder~\cite{who2022wmhr}. More than 720{,}000 people die by suicide
each year, and suicide is the third leading cause of death among
individuals aged 15--29~\cite{who2024suicide}. Depression and anxiety
are also estimated to cost the global economy approximately 12 billion
working days annually~\cite{who2022work}. Because early intervention
substantially improves outcomes, detection at the earliest feasible
stage is a public-health priority.

Current approaches fall short in three distinct ways. First, clinical
instruments such as structured interviews and validated questionnaires
are accurate but expensive, require trained administration, and depend
entirely on what the respondent is willing or able to report. Second,
most computational systems rely on a single modality---text, speech, or
facial expression---even though the relevant signal is distributed
across modalities rather than localised in any one of them. Third, and
most importantly for this work, existing systems do not explicitly
model \emph{disagreement} between signals. When a respondent says ``I am
fine'' while their language becomes more absolutist, their speech slows,
and their facial expressivity flattens, the conflict itself carries
information. Conventional fusion architectures treat such conflict as
noise to be averaged away.

We distinguish two kinds of evidence. \emph{Explicit} evidence is what
the individual consciously and intentionally reports: ordinal responses
to structured questions gathered through a branching conversational
interface. \emph{Implicit} evidence is what the system infers from
behaviour---linguistic markers, response latency, prosody, and facial
dynamics---and is more difficult to deliberately control. We define the
\emph{consistency score} $C$ as the discrepancy between the two
representations. A large $C$ indicates that the individual's explicit
account and implicit behaviour diverge, which is consistent with social
desirability bias, emotional masking, or limited insight.

This paper makes three design contributions.
\begin{enumerate}
  \item \textbf{A dual-elicitation conversational architecture.} An
        FSM-based interview engine with eight clinical modules, grouped
        into five thematic domains, designed to collect explicit
        self-report and implicit multimodal evidence within a single
        structured interaction.
  \item \textbf{Consistency modeling as a first-class signal.} A
        projection maps implicit features into the explicit feature
        space, and the residual is quantified as $C$. We separate the
        \emph{dual role} of $C$---as a potential label-informative
        signal and as a hypothesised error-predicting reliability
        index---and argue that these roles must not be conflated in a
        single additive risk term.
  \item \textbf{Population-relative statistical calibration.} A
        baseline layer converts absolute risk and latency into
        population-relative z-scores, producing interpretable anomaly
        flags that would route ambiguous or atypical sessions to human
        review.
\end{enumerate}

\subsection{Scope and Status of the Claims}
\label{sec:scope}

This is a \emph{conceptual framework paper}. It presents an
architecture, a formal notation, a set of design arguments, and a
falsifiable evaluation plan. \textbf{It reports no experimental
results.} The speech, facial, and behavioural components are specified
in full and grounded in public benchmarks, but no joint evaluation has
been carried out. Wherever an empirical claim would be needed, this
paper states a hypothesis and identifies the experiment that would test
it.

We state this up front deliberately. A framework paper is a legitimate
genre, but it is a different contribution from an empirical study, and
its claims must be assessed accordingly. Readers seeking validated
performance figures should treat this work as a specification and
protocol, not as evidence that the system works.

% =====================================================================
\section{Related Work}

\subsection{Text-Based Approaches}

Coppersmith et al.~\cite{coppersmith2015mental} demonstrated that
self-disclosed mental-health diagnoses on Twitter can be classified from
language using NLP features and supervised models. Their work established
that linguistic markers carry detectable signal, but it remained
text-only and relied on voluntary public disclosure of diagnosis.
General-purpose contextual encoders such as BERT~\cite{devlin2019bert}
and RoBERTa~\cite{liu2019roberta} substantially improved text
representation learning and are used here as \emph{encoders}, not as
mental-health models in their own right; neither incorporates
non-linguistic signal. DistilBERT~\cite{sanh2019distilbert} retains most
of BERT's task performance at a fraction of the inference cost and is
the encoder proposed for this framework.

\subsection{Speech-Based Analysis}

The Distress Analysis Interview Corpus (DAIC-WOZ) was introduced by
Gratch et al.~\cite{gratch2014daic} as a corpus of clinical interviews
with associated depression severity labels. Alhanai et
al.~\cite{alhanai2018daic} subsequently showed that audio and text
sequence models over DAIC-WOZ can detect depression, but the approach
did not integrate facial or behavioural evidence. We cite the corpus and
the modelling work separately to avoid attributing the dataset to the
modelling study.

\subsection{Physiological and EEG-Based Systems}

Picard~\cite{picard1997affective} established affective computing as a
field, arguing that affect should be treated as an input signal amenable
to computational modelling. Koelstra et al.~\cite{koelstra2012deap}
released DEAP, a widely used physiological benchmark for emotion
analysis that has supported SVM, random forest, and CNN baselines.
Physiological sensing nonetheless requires specialised hardware (EEG
headsets, wearable sensors), which limits deployment for low-cost
screening. Russell~\cite{russell1980circumplex} proposed the
circumplex model, representing affect on continuous valence and arousal
axes rather than as discrete categories; this representation underlies
the facial component of the present framework.

\subsection{Facial Emotion Recognition}

Goodfellow et al.~\cite{goodfellow2013fer} released FER2013, an early
large-scale benchmark of 35{,}887 grayscale facial images annotated with
seven emotion classes. Mollahosseini et
al.~\cite{mollahosseini2017affectnet} introduced AffectNet, which added
continuous valence and arousal annotations for in-the-wild images and
therefore supports dimensional rather than purely categorical modelling.
Facial action units, grounded in the Facial Action Coding
System~\cite{ekman1978facs}, provide an alternative descriptive
vocabulary. Kossaifi et al.~\cite{kossaifi2019action} explored action
unit representations on AU-annotated datasets such as BP4D and DISFA,
producing richer descriptors than discrete labels; that line of work
remains facial-only and does not incorporate speech or behaviour.

\subsection{Multimodal Fusion}

Baltru\v{s}aitis et al.~\cite{baltrusaitis2018multimodal} provided the
taxonomy of multimodal machine learning---representation, translation,
alignment, fusion, and co-learning---that organises much of the
subsequent literature. Tensor Fusion Network~\cite{zadeh2017tensor} is a
representative fusion architecture: it models unimodal, bimodal, and
trimodal interactions through an outer product of modality embeddings.
Such architectures assume that modalities are aligned and mutually
reinforcing; disagreement between modalities is treated as noise rather
than as a measurable quantity. MODMA~\cite{liu2022modma} and
related efforts supply clinically collected multimodal data for
depression analysis, while Xu et al.~\cite{xu2025multimodal} fuse BERT
text features with Wav2Vec~2.0 speech features for depression
prediction; their contribution is a better fusion layer, not a model of
agreement or disagreement.

\subsection{Consistency, Response Bias, and Masking}

The discrepancy between self-report and inferred state has a long
history in psychological measurement. Social desirability
bias---the tendency to present oneself favourably---was formalised by
Crowne and Marlowe~\cite{crowne1960social} and remains a documented
source of error in self-report instruments. In the mental-health
context, the same phenomenon appears as emotional masking or
minimisation. Computational systems typically either ignore this
discrepancy or treat it as label noise. We are not aware of a widely
adopted framework that quantifies explicit--implicit inconsistency as
an explicit, inspectable feature and then uses it to modulate
confidence.

\subsection{Research Gap}
\label{sec:gap}

Three gaps motivate this work. (i) \emph{Inconsistency is not modeled.}
Fusion architectures optimise agreement; they do not represent
disagreement as an outcome of interest. (ii) \emph{Global coordination
is absent.} Single-modality and simple fusion models still lack a
mechanism that reconciles evidence across modalities at a session
level. (iii) \emph{Statistical grounding is missing.} Most systems
report an absolute score with no reference to a population distribution,
so a given score cannot be interpreted as typical or atypical for a
comparable individual. Figure~\ref{fig:research_gap} summarises these
gaps against the reviewed literature.

\begin{figure}[H]
    \centering
    \includegraphics[width=\textwidth]{new.png}
    \caption{Research gap analysis across mental-health detection
    approaches. Text-only, speech-only, physiological, and facial
    systems each leave inconsistency modeling and population-relative
    calibration unaddressed; fusion systems optimise agreement rather
    than representing disagreement.}
    \label{fig:research_gap}
\end{figure}

% =====================================================================
\section{Methodology}

\subsection{System Overview}

Figure~\ref{fig:architecture} shows the proposed end-to-end data flow. A
user completes the FSM interview three times. Repeated administration is
deliberate: a single session may reflect a transient state, whereas
patterns that recur across sessions are more likely to be stable and
therefore more trustworthy as evidence. Each session would yield a set
of modality-specific feature vectors, a consistency score, and a
calibrated risk estimate.

\begin{figure}[H]
\centering
\includegraphics[width=0.95\linewidth]{dg.png}
\caption{Proposed system architecture. The FSM interviewing engine
collects explicit self-reports and implicit multimodal evidence;
features are fused into a joint vector; the consistency score $C$ is
computed in feature space; population-relative z-score calibration
situates the session against a reference distribution; and the
reliability gate routes high-$C$ or statistically anomalous sessions to
human review.}
\label{fig:architecture}
\end{figure}

\subsection{Data Sources and Evaluation Cohorts}
\label{sec:data}

The framework draws on two qualitatively different kinds of data, and we
keep their roles strictly separate.

\paragraph{Target deployment data.} In deployment, all modalities would
be collected in a single consented session: ordinal self-reports,
free-text responses, audio, and facial video, together with behavioural
logs from the FSM engine.

\paragraph{Component grounding benchmarks.} No single public dataset
provides all four modalities under one protocol. The speech component is
therefore specified against DAIC-WOZ~\cite{gratch2014daic}, which
supplies clinical interview audio and transcripts with depression
severity labels, and MODMA~\cite{liu2022modma}, which supplies
clinically collected voice recordings with emotional-state annotations.
The publicly distributed DAIC-WOZ release provides audio and transcripts
rather than raw facial video; consequently the facial component is
grounded separately in FER2013~\cite{goodfellow2013fer} and
AffectNet~\cite{mollahosseini2017affectnet}. AffectNet provides over
450{,}000 in-the-wild images annotated with both discrete emotion labels
and continuous valence--arousal values, which is what makes dimensional
modelling feasible. Text representations can be trained or fine-tuned on
the transcripts available in DAIC-WOZ and MODMA.

These benchmarks inform component design and extractor selection. They
do not constitute a joint multimodal evaluation, and the absence of a
single corpus containing all modalities is itself a constraint on the
evaluation plan (Section~\ref{sec:evaluation}).

\paragraph{Behavioural data.} Per-question response latency, session
duration, interaction count, and ordinal self-report values would be
logged by the FSM engine itself.

\subsection{Dual Elicitation Design}

\subsubsection{The Interview Engine}

The engine is a finite state machine
\begin{equation}
\pi = (s_0,\; s_1,\; \ldots,\; s_T), \qquad s_T \in F,
\end{equation}
where each transition is conditioned on the preceding response. The
interview comprises \textbf{eight clinical modules} organised into
\textbf{five thematic domains} (Table~\ref{tab:modules}). Branching
logic concentrates follow-up questions where the preceding answer
indicates elevated concern; for example, a response indicating marked
social isolation triggers additional loneliness-focused items. The
collection of modalities is
\begin{equation}
S = S_{\text{text}} \cup S_{\text{speech}} \cup S_{\text{behaviour}}
    \cup S_{\text{facial}} \cup S_{\text{consistency}}.
\end{equation}

\begin{table}[H]
\centering
\caption{Interview structure: eight clinical modules across five
thematic domains.}
\label{tab:modules}
\begin{tabular}{@{}ll@{}}
\toprule
Thematic domain & Clinical modules \\
\midrule
Physical well-being & Sleep; health \\
Mood and stress     & Mood; stress \\
Social and self     & Social; self \\
Functioning         & Study \\
Safety              & Safety \\
\bottomrule
\end{tabular}
\end{table}

\subsubsection{Phase 1: Structured Elicitation (Explicit Features)}

\begin{equation}
\mathbf{x}^{\text{explicit}} = \bigl[\, x_{\text{sleep}},\; x_{\text{mood}},\;
x_{\text{stress}},\; x_{\text{social}},\; x_{\text{self}},\;
x_{\text{study}},\; x_{\text{health}},\; x_{\text{safety}} \,\bigr]^{\top}
\in \mathbb{R}^{m}, \quad m = 8.
\end{equation}

\noindent These values encode what the user \emph{consciously reports} via
normalised ordinal scale responses.

\subsubsection{Phase 2: Conversational Elicitation (Implicit Features)}

\begin{equation}
\mathbf{x}^{\text{implicit}} = \bigl[\, \mathbf{e}_{\text{embed}},\;
x_{\text{sent}},\; x_{\text{self\_ref}},\; x_{\text{dist}},\;
x_{\text{intensity}} \,\bigr]^{\top} \in \mathbb{R}^{k},
\end{equation}

\noindent where $\mathbf{e}_{\text{embed}} \in \mathbb{R}^{768}$ and
$k = 772$. These values encode what the user's language and behaviour
\emph{implicitly reveal}, independent of conscious self-report.
Table~\ref{tab:vectors} summarises the modality vectors and their
dimensionality.

\begin{table}[H]
\centering
\caption{Modality-specific feature vectors. Speech dimensionality depends
on the selected openSMILE feature set.}
\label{tab:vectors}
\begin{tabular}{@{}llc@{}}
\toprule
Modality & Features & Dimension \\
\midrule
Explicit self-report & Eight ordinal clinical dimensions & 8 \\
Implicit text & DistilBERT CLS, VADER, self-reference, distortion, intensity & 772 \\
Facial & $\bar{V}$, $\bar{A}$, $\bar{D}$, EAR, MAR & 5 \\
Behavioural & Mean latency, latency SD, session duration, questions answered & 4 \\
Speech & MFCCs, pitch, jitter, shimmer, energy & variable \\
Consistency & Score $C$ and its cosine variant $C_{\cos}$ & 2 \\
\bottomrule
\end{tabular}
\end{table}

\subsection{Consistency Modeling}
\label{sec:consistency}

\subsubsection{Learned Projection}

A linear map $\phi:\mathbb{R}^{k}\rightarrow\mathbb{R}^{m}$ projects
implicit features into the explicit feature space:
\begin{equation}
\hat{\mathbf{x}}^{\text{explicit}}
= \phi\bigl(\mathbf{x}^{\text{implicit}}\bigr)
= \mathbf{W}\,\mathbf{x}^{\text{implicit}} + \mathbf{b},
\qquad \mathbf{W}\in\mathbb{R}^{m\times k},\; \mathbf{b}\in\mathbb{R}^{m},
\end{equation}
trained by least squares:
\begin{equation}
\mathcal{L}_{\phi}
= \left\| \mathbf{x}^{\text{explicit}}
        - \phi\bigl(\mathbf{x}^{\text{implicit}}\bigr) \right\|_{2}^{2}.
\end{equation}

In an initial implementation, $\phi$ could be initialised heuristically
before being fitted. For example, the mood dimension might be estimated
as
\begin{equation}
\hat{x}_{\text{mood}}
= w_{1}\,x_{\text{sent}} + w_{2}\,x_{\text{self\_ref}}
+ w_{3}\,x_{\text{dist}} + b .
\end{equation}

\subsubsection{Inconsistency Score}

\begin{equation}
C = \left\| \mathbf{x}^{\text{explicit}}
        - \hat{\mathbf{x}}^{\text{explicit}} \right\|_{2}
= \sqrt{\sum_{i=1}^{m}\bigl(x_i^{\text{explicit}}
        - \hat{x}_i^{\text{explicit}}\bigr)^{2}} .
\end{equation}

\noindent A cosine-based variant that is insensitive to vector magnitude is
\begin{equation}
C_{\cos} = 1
- \frac{\mathbf{x}^{\text{explicit}} \cdot \hat{\mathbf{x}}^{\text{explicit}}}
       {\left\|\mathbf{x}^{\text{explicit}}\right\|_{2}\,
        \left\|\hat{\mathbf{x}}^{\text{explicit}}\right\|_{2}} .
\end{equation}

A large $C$ indicates that linguistic responses contradict numerical
self-reports, which is consistent with emotional masking or social
desirability bias~\cite{crowne1960social}.

\subsection{Feature Extraction}

\subsubsection{Text Features}

\emph{Semantic embedding.} DistilBERT~\cite{sanh2019distilbert} provides
a six-fold inference speed advantage over BERT while retaining most of
its task performance:
\begin{equation}
\mathbf{e}_{\text{embed}} = \text{DistilBERT}(\text{response})[\texttt{CLS}]
\in \mathbb{R}^{768}.
\end{equation}

\emph{Sentiment polarity.} VADER~\cite{hutto2014vader} returns a
compound polarity score:
\begin{equation}
x_{\text{sent}} = \text{VADER}(\text{response}) \in [-1,\,1].
\end{equation}

\emph{Self-reference ratio.} Elevated first-person singular usage is an
established linguistic marker of depressive
rumination~\cite{coppersmith2015mental}:
\begin{equation}
x_{\text{self\_ref}}
= \frac{\text{count(first-person singular pronouns)}}
       {N_{\text{words}}} .
\end{equation}

\emph{Cognitive distortion score.} Absolutist terms such as ``always'',
``never'', ``nothing'', and ``completely'' are detected via a predefined
lexicon:
\begin{equation}
x_{\text{dist}}
= \frac{\text{count(distortion keywords)}}{N_{\text{words}}} .
\end{equation}

\noindent The implicit vector is therefore
$\mathbf{x}^{\text{implicit}}\in\mathbb{R}^{772}$. The distortion and
self-reference terms capture clinical signal that sentiment polarity
alone does not.

\subsubsection{Facial Features via Dimensional Emotion Modeling}

Discrete categorical labels cannot represent continuously graded
affective states, which limits their usefulness for screening. Following
Russell's circumplex model~\cite{russell1980circumplex} as operationalised
in AffectNet~\cite{mollahosseini2017affectnet}, the facial component maps
CNN emotion probabilities onto continuous valence and arousal and
derives a distress index.

\emph{Emotion probability distribution.} A CNN outputs a softmax
distribution over $K=8$ categories:
\begin{equation}
P(\text{emotion}_i) = \frac{e^{z_i}}{\sum_{j=1}^{K} e^{z_j}},
\qquad \sum_{i=1}^{K} P(\text{emotion}_i) = 1,
\end{equation}
where $i \in \{\text{happy, sad, anger, fear, surprise, disgust, neutral,
contempt}\}$.

\emph{Valence.} Valence $V\in[-1,1]$ encodes the positivity of the
expressed state:
\begin{equation}
V = \sum_{i=1}^{K} P(\text{emotion}_i)\,w_V(\text{emotion}_i).
\end{equation}

\emph{Arousal.} Arousal $A\in[-1,1]$ encodes activation intensity:
\begin{equation}
A = \sum_{i=1}^{K} P(\text{emotion}_i)\,w_A(\text{emotion}_i).
\end{equation}

\noindent The weights $w_V$ and $w_A$ are \textbf{heuristic priors}
chosen to respect the circumplex ordering of each category; they are
not regression coefficients estimated from AffectNet's continuous
annotations. The values proposed for the initial implementation are:
$w_V(\text{happy})=+1.0$, $w_V(\text{surprise})=+0.4$,
$w_V(\text{neutral})=0.0$, $w_V(\text{contempt})=-0.7$,
$w_V(\text{sad})=-0.7$, $w_V(\text{fear})=-0.8$,
$w_V(\text{anger})=-0.9$, $w_V(\text{disgust})=-0.95$; and
$w_A(\text{fear})=+0.9$, $w_A(\text{anger})=+0.8$,
$w_A(\text{surprise})=+0.7$, $w_A(\text{happy})=+0.6$,
$w_A(\text{contempt})=+0.3$, $w_A(\text{neutral})=0.0$,
$w_A(\text{disgust})=-0.2$, $w_A(\text{sad})=-0.3$. These priors are a
design choice, and Section~\ref{sec:limitations} records them as a
limitation.

\emph{Distress index.} A per-frame distress score combines negative
valence with activation magnitude:
\begin{equation}
D = (1 - V)\,|A| .
\end{equation}

\noindent For fear ($V=-0.8$, $A=0.9$), $D = 1.62$; for anger
($V=-0.9$, $A=0.8$), $D = 1.52$; for sadness ($V=-0.7$, $A=-0.3$),
$D = 0.51$; for happiness ($V=+1.0$, $A=+0.6$) and neutrality
($V=0$, $A=0$), $D = 0$.

This formulation has an important limitation that we state explicitly:
$D$ collapses to zero both for genuinely positive affect and for flat
or blunted affect. It therefore \emph{cannot} by itself distinguish a
calm individual from one exhibiting emotional blunting. We do not claim
otherwise. A candidate discriminator---session-level expressivity
variance, $\sigma_V$ and $\sigma_A$ computed over the session---is
identified in Section~\ref{sec:future} but is not validated here. The
distress index should be read as a measure of negative high-arousal
affect, not as a measure of flatness.

\emph{Temporal smoothing.} Single-frame predictions are noisy; a sliding
window of $N$ frames is maintained:
\begin{equation}
\bar{V} = \frac{1}{N}\sum_{n=1}^{N}V_n, \qquad
\bar{A} = \frac{1}{N}\sum_{n=1}^{N}A_n, \qquad
\bar{D} = \frac{1}{N}\sum_{n=1}^{N}D_n .
\end{equation}

\noindent A default window of $N=30$ frames (approximately one second at
30\,FPS) is proposed to suppress single-frame misclassification while
preserving genuine transitions.

\emph{Micro-behavioural features.} MediaPipe
FaceMesh~\cite{lugaresi2019mediapipe} provides facial landmarks from
which the eye aspect ratio (EAR) and mouth aspect ratio (MAR) are
computed:
\begin{equation}
\text{EAR} = \frac{\lVert p_{\text{upper}} - p_{\text{lower}}\rVert_{2}}
                 {\lVert p_{\text{left}} - p_{\text{right}}\rVert_{2}},
\qquad
\text{MAR} = \frac{d_{\text{upper--lower lip}}}
                  {d_{\text{mouth corner distance}}} .
\end{equation}

\emph{Session-level facial vector.}
\begin{equation}
\mathbf{x}^{\text{face}} = \bigl[\,\bar{V},\; \bar{A},\; \bar{D},\;
\text{EAR},\; \text{MAR}\,\bigr]^{\top} \in \mathbb{R}^{5}.
\end{equation}

\subsubsection{Speech Features}

Speech would be analysed with librosa~\cite{mcfee2015librosa} and
openSMILE~\cite{eyben2010opensmile} to obtain MFCCs, pitch, jitter,
shimmer, and energy. These prosodic and spectral features enter the
implicit space. Feature selection is informed by the acoustic analyses
reported on DAIC-WOZ~\cite{gratch2014daic,alhanai2018daic} and
MODMA~\cite{liu2022modma}.

\subsubsection{Behavioural Features}

\begin{equation}
\mathbf{x}^{\text{behavior}}
= \bigl[\,\bar{L},\; \sigma^{\text{sess}}_{L},\; D_{s},\; Q\,\bigr]^{\top}
\in \mathbb{R}^{4},
\end{equation}
where
\begin{equation}
\bar{L} = \frac{1}{T}\sum_{t=1}^{T}L_t,
\qquad
\sigma^{\text{sess}}_{L}
= \sqrt{\frac{1}{T}\sum_{t=1}^{T}\bigl(L_t - \bar{L}\bigr)^{2}} ,
\end{equation}
$L_t$ is the response latency for question $t$,
$D_{s}$ is total session duration, and $Q$ is the number of questions
answered. We use $L$ for latency throughout and reserve $\lambda$ for
no other purpose. All features are standardised prior to fusion.

\subsection{Feature Fusion and Risk Estimation}

\subsubsection{Early Fusion}

All modality vectors are concatenated before risk estimation:
\begin{equation}
\mathbf{x}^{\text{fused}}
= \bigl[\,\mathbf{x}^{\text{explicit}},\; \mathbf{x}^{\text{implicit}},\;
\mathbf{x}^{\text{face}},\; \mathbf{x}^{\text{speech}},\;
\mathbf{x}^{\text{behavior}}\,\bigr]^{\top}.
\end{equation}

\noindent Because $C$ is a function of the same features, including it
here would duplicate information and complicate the interpretation of
the gate. It is therefore computed separately and excluded from
$\mathbf{x}^{\text{fused}}$.

\subsubsection{Base Risk Estimate}

\emph{Initial phase (weighted linear combination).}
\begin{equation}
r = w_{1}\,\bar{x}^{\text{explicit}}
  + w_{2}\,\bar{x}^{\text{behavior}}
  + w_{3}\,\bar{D},
\qquad
\bar{x}^{\text{explicit}} = \frac{1}{m}\sum_{i=1}^{m}x_i^{\text{explicit}},
\label{eq:linear-risk}
\end{equation}
with $w_{1}=0.55$, $w_{2}=0.25$, $w_{3}=0.20$. These weights are
starting heuristics and are replaced by a learned model in the second
phase.

\emph{Second phase (learned risk models).} XGBoost~\cite{chen2016xgboost}
is proposed as the primary learner; random forest as the baseline, and
LightGBM as an alternative comparison. The learned estimate is
$\hat{r} = f(\mathbf{x}^{\text{fused}})$, and a PCA projection
$\mathbf{z} = \text{PCA}(\mathbf{x}^{\text{fused}})$ is proposed for
diagnostic visualisation.

\subsubsection{Decision Rule}

The final decision combines the point estimate, the reliability index,
and the statistical baseline (Section~\ref{sec:baseline}):
\begin{equation}
\text{escalate}(i) =
\mathbb{1}\bigl[\,C_i > \tau_C \;\lor\; |Z_{R,i}| > 2
\;\lor\; |Z_{L,i}| > 2 \;\lor\; \text{safety}_i = 1\,\bigr].
\label{eq:decision}
\end{equation}

\noindent The reported output would be the pair $\bigl(\hat{r}_i,
\hat{\rho}(C_i)\bigr)$ together with the escalation flag. The point
estimate is never inflated by an inconsistency penalty.

\subsubsection{Planned Extension: Prediction-Level Divergence}
\label{sec:divergence}

The present consistency score operates in \emph{feature space}. A natural
extension assigns each modality its own risk estimate and measures
disagreement in \emph{prediction space}: $D_{ET}$ between explicit and
text-implicit predictions, $D_{EF}$ between explicit and facial
predictions, and $D_{TF}$ between text and facial predictions. This
would allow conflict to be localised to specific modalities rather than
reported as a single scalar. It is not developed here.

\subsection{Statistical Baseline Modeling and Anomaly Detection}
\label{sec:baseline}

\subsubsection{Population Distribution}

For $N$ reference sessions with risk estimates $R_i$ and mean response
latencies $\bar{L}_i$:
\begin{equation}
\mu_R = \frac{1}{N}\sum_{i=1}^{N}R_i,
\qquad
\sigma_R = \sqrt{\frac{1}{N}\sum_{i=1}^{N}\bigl(R_i - \mu_R\bigr)^{2}},
\end{equation}
\begin{equation}
\mu_L = \frac{1}{N}\sum_{i=1}^{N}\bar{L}_i,
\qquad
\sigma^{\text{pop}}_{L}
= \sqrt{\frac{1}{N}\sum_{i=1}^{N}\bigl(\bar{L}_i - \mu_L\bigr)^{2}} .
\end{equation}

\noindent The reference distribution would be estimated empirically from
a reference cohort. No synthetic population is assumed; a small
reference cohort would yield unstable $\mu$ and $\sigma$ estimates, which
Section~\ref{sec:limitations} records as a limitation.

\subsubsection{Z-Score Normalization}

\begin{equation}
Z_R = \frac{R_i - \mu_R}{\sigma_R},
\qquad
Z_L = \frac{\bar{L}_i - \mu_L}{\sigma^{\text{pop}}_{L}} .
\end{equation}

\noindent Z-scores provide the population-relative context that absolute
scores lack: a given score is interpretable only against what is typical
for comparable sessions.

\subsubsection{Towards Personalised Baselines}

A population baseline cannot distinguish between-subject from
within-subject variation. An individual who is habitually quiet may be
entirely typical for themselves while deviating from the population. The
planned blend is
\begin{equation}
S_i = \kappa_i\,\frac{R_i - \mu_R}{\sigma_R}
    + (1-\kappa_i)\,\frac{R_i - \mu_i}{\sigma_i},
\end{equation}
where $\mu_i$ and $\sigma_i$ would be estimated from the individual's
own longitudinal sessions and $\kappa_i \in [0,1]$ controls the blend.
$\kappa_i$ is \emph{not} set in this work; the formulation is included to
fix notation for the planned extension. Note that $\kappa$ is used here
and nowhere else.

\subsubsection{Anomaly Flags}

A session would be flagged when a z-score exceeds two standard
deviations, or when consistency is high, or when the safety module is
positive (Eq.~\ref{eq:decision}). Flagged sessions would be routed to a
clinically trained reviewer; they would never be auto-diagnosed.

% =====================================================================
\section{Design Rationale: Consistency as Reliability, Not Severity}
\label{sec:dualrole}

The most consequential design decision in this framework concerns how
$C$ enters the risk estimate. We argue for treating $C$ as a
\emph{reliability} signal rather than as an additive severity term. The
argument has three parts.

\paragraph{Discrepancy has no inherent clinical direction.} $C$ is an
unsigned distance between two representations. It increases when a
respondent under-reports distress while their language suggests
otherwise---a masking pattern. But it also increases when the
projection $\phi$ is poorly estimated, when a question is ambiguous,
when speech extraction is corrupted by noise, or when a respondent
simply answers carelessly. The same scalar therefore mixes clinically
meaningful incongruence with measurement error. Mapping $C$ directly
onto risk requires an assumption that has to be justified, not assumed.

\paragraph{Adding $C$ to risk double-counts it.} If $C$ is also present
among the fusion features---as it would be under a naive
implementation---then including it again as an additive term counts the
same evidence twice. Even if $C$ is excluded from the feature vector,
an additive term imposes a monotone, unbounded severity interpretation
on a quantity whose scale depends on the dimensionality and
normalisation of the explicit space.

\paragraph{A reliability reading is testable and safer.} The framework
therefore treats $C$ as a hypothesised predictor of estimation error.
The reliability index is defined as the local accuracy in a
neighbourhood $\mathcal{B}(C)$ of the consistency score:
\begin{equation}
\hat{\rho}(C)
= \frac{1}{|\mathcal{B}(C)|}
  \sum_{j\in\mathcal{B}(C)} \mathbb{1}\bigl[\hat{y}_j = y_j\bigr] .
\end{equation}

\noindent A session is routed to review when
\begin{equation}
\text{review}(i) = \mathbb{1}\bigl[\,C_i > \tau_{C}\,\bigr],
\label{eq:gate}
\end{equation}
where $\tau_C$ would be selected on held-out data to trade coverage
against precision. This preserves the clinically important case in which
a respondent reports low distress while implicit signals disagree: such
a session is escalated for human assessment rather than assigned an
artificially high probability by an additive penalty.

The design makes a strong, falsifiable commitment: if $C$ does not
predict error (Hypothesis H2, Section~\ref{sec:evaluation}), the
reliability-gating rationale fails and the mechanism should be
redesigned or discarded. That is preferable to an additive formulation
whose contribution cannot be falsified because it always increases the
score.

% =====================================================================
\section{Proposed Evaluation Protocol}
\label{sec:evaluation}

Because this paper reports no results, we specify the study that would
test the framework. The protocol is intended to be pre-registerable.

\subsection{Hypotheses}

\begin{enumerate}
  \item[\textbf{H1}] \emph{Label association.} The distribution of $C$
        differs between individuals with and without a depression
        indication, with higher $C$ more frequent in the depressed
        group.
  \item[\textbf{H2}] \emph{Reliability.} Classification accuracy
        decreases as $C$ increases, and $C$ predicts per-sample error
        beyond what class balance explains.
  \item[\textbf{H3}] \emph{Gating utility.} Gating on $C$ increases
        precision within the retained subset at a quantifiable cost in
        coverage, relative to ungated predictions.
  \item[\textbf{H4}] \emph{Multimodal benefit.} The fused multimodal
        model outperforms the best single-modality baseline on held-out
        data.
\end{enumerate}

H2 is the load-bearing hypothesis: it is the empirical foundation of the
design decision argued in Section~\ref{sec:dualrole}. H4 is the claim
most often asserted without support in this literature, and it must be
tested rather than assumed.

\subsection{Data Requirements}

No single public dataset contains all four modalities under one
protocol. Two designs are therefore possible. \emph{Design A} evaluates
each component on its native benchmark and reports component-level
results only; it cannot test H4. \emph{Design B} collects a new consented
cohort with a common protocol, which is required to test H4 properly.
We recommend Design B and treat Design A as an interim step.

\subsection{Splitting and Statistics}

Splits must be \emph{subject-independent}: the same participant must not
appear in both training and test partitions, since repeated sessions
from one individual would otherwise leak identity and inflate accuracy.
We propose repeated stratified $k$-fold cross-validation with multiple
random seeds, reporting mean $\pm$ standard deviation and confidence
intervals. Because seed-to-seed variation is itself an estimate,
inferential claims should use paired tests across folds with effect
sizes and correction for multiple comparisons. A small number of seeds
should be reported as exploratory, not confirmatory.

\subsection{Metrics}

Classification performance would be reported as accuracy, precision,
recall, $F_1$, and area under the ROC curve, with calibration assessed
by expected calibration error---since the framework claims calibrated
output, discrimination alone is insufficient. Facial valence and arousal
would be evaluated with the concordance correlation coefficient used in
the AffectNet protocol~\cite{mollahosseini2017affectnet}:
\begin{equation}
\rho_c =
\frac{2\rho\,\sigma_{\hat{\theta}}\,\sigma_{\theta}}
     {\sigma_{\hat{\theta}}^{2} + \sigma_{\theta}^{2}
      + (\mu_{\hat{\theta}} - \mu_{\theta})^{2}} .
\end{equation}
Severity association would use the Spearman rank correlation:
\begin{equation}
\rho_s\bigl(C,\,y\bigr)
= \frac{\operatorname{Cov}\bigl(\mathrm{rank}(C),\,\mathrm{rank}(y)\bigr)}
       {\sigma_{\mathrm{rank}(C)}\,\sigma_{\mathrm{rank}(y)}} .
\end{equation}

\subsection{Ablations}

Four configurations would isolate the consistency mechanism:
(i)~explicit self-report only; (ii)~explicit plus implicit linguistic
features; (iii)~configuration (ii) plus behavioural features;
(iv)~configuration (iii) with $C$ used as a reliability gate. A fifth
configuration using $C$ as an additive penalty would be included
specifically to test the design argument of Section~\ref{sec:dualrole}.
Each configuration would report accuracy, $F_1$, and coverage.

\subsection{Clinical Ground Truth}

Screening performance should ultimately be calibrated against
established instruments---PHQ-9~\cite{kroenke2001phq9} and
GAD-7~\cite{spitzer2006gad7}---with ROC analysis and threshold selection
justified on clinical rather than statistical grounds. Agreement with an
instrument is not the same as clinical validity, and neither substitutes
for diagnostic assessment.

\subsection{Threats to Validity}

\emph{Dataset shift}: models trained on benchmark corpora may not
transfer to the deployment population. \emph{Label quality}: public
Reddit-based labels are self-disclosed and unverified.
\emph{Demographic bias}: facial and prosodic models are known to vary
across groups, and the affect priors derive from a predominantly Western
annotation protocol. \emph{Confounding}: latency and expressivity are
affected by device, bandwidth, and environment, not only by affect.
\emph{Construct validity}: $C$ measures representational disagreement,
which is not the same construct as masking; H1 and H2 test association,
not the mechanism.

% =====================================================================
\section{Discussion}

\subsection{What the Framework Contributes}

The framework's contribution is conceptual rather than empirical. It
does three things that current fusion architectures do not. It makes
explicit--implicit inconsistency an inspectable quantity rather than an
implicit assumption of agreement. It separates the question ``how
severe is this case?'' from ``how much should we trust this estimate?'',
which are routinely conflated when a single confidence-like score is
folded into a prediction. And it places an absolute score in a
population context, so that a result can be described as typical or
atypical rather than merely high or low.

\subsection{Dimensional Facial Modeling}

The facial component departs from categorical classification by
representing affect on continuous valence and arousal axes, motivated by
the circumplex model~\cite{russell1980circumplex} and by AffectNet's
annotation protocol~\cite{mollahosseini2017affectnet}.

The claim that dimensional modelling distinguishes calm from emotional
blunting is \emph{not} supported by the present formulation. Both
states map near $V \approx 0$, $A \approx 0$, and therefore both yield
$D \approx 0$. What the model captures is negative high-arousal affect
(fear, anger), and this is what the distress index should be understood
to measure. Distinguishing blunting would require an additional measure
of expressivity dynamics, such as session-level $\sigma_V$ and
$\sigma_A$, which is proposed but untested.

Similarly, the weights $w_V$ and $w_A$ are heuristic priors that respect
the circumplex ordering of the eight categories. They are not derived
from AffectNet's continuous annotations and should not be described as
such. Refitting them by regression against the dimensional targets, and
evaluating with the concordance correlation coefficient, is future work.

\subsection{Why a Framework Paper Is Appropriate Here}

Multimodal mental-health detection has a replication problem: fusion
results are frequently reported on benchmarks that do not share a
protocol with the intended deployment setting, and the architectural
claims---that more modalities are better, that disagreement is noise---are
rarely stated precisely enough to be falsified. Specifying the
architecture and the test before running it is a way of avoiding that
failure mode. The cost is that this paper provides no evidence that the
system works. That cost is real and should not be minimised: a framework
is a proposal, and proposals can be wrong.

% =====================================================================
\section{Ethical Considerations}

\subsection{Data Governance}

This paper proposes a framework and reports no data collection. Any
future study would require independent ethics approval, informed
consent, and a data-protection agreement. Use of public benchmarks
would be governed by the consent and ethics approvals obtained by the
original data providers.

\subsection{Privacy and Scope of Inference}

Sensitive personal data would be anonymised and encrypted. The system
is not proposed as an automated diagnostic device. The consistency
score is intended as an internal quantity: it would not be surfaced to
users in raw form, and it would not be presented as a clinical finding.
Anomaly flags would be directed exclusively to clinically trained
reviewers. Bias monitoring and fairness evaluation across demographic
groups are prerequisites for deployment, not optional extras.

\subsection{Safety Escalation}
\label{sec:safety}

The interview includes a safety module, and the explicit feature vector
contains a safety dimension. The system must therefore be designed on
the assumption that suicidal ideation will be disclosed.

Anomaly flags for \emph{asynchronous} clinical review are not a
sufficient response to acute risk. We specify a separate, immediate
pathway. If the safety module records a positive response, the
screening flow halts, crisis-contact information and local emergency
service numbers are displayed in-session, and the designated on-call
clinician is notified according to a pre-agreed escalation protocol
with a defined response window. The system does not provide emergency
response itself and is not a crisis service; this must be stated to
users before participation, during informed consent, and on every
escalation screen.

A deployed version handling disclosures of this kind would require
clinical governance, an on-call rota, and documented follow-up, none of
which is evaluated in this paper.

\subsection{Transparency}

Model behaviour, the reliability index, and the escalation rule are
specified so that they can be audited. The heuristic status of the
valence and arousal priors is disclosed (Section~\ref{sec:limitations}).

% =====================================================================
\section{Limitations}
\label{sec:limitations}

\emph{No empirical evaluation.} This paper reports no results. Every
performance-related claim in the framework is a hypothesis, not a
finding.

\emph{No single multimodal corpus.} No public dataset contains text,
speech, face, and behaviour under one protocol, so the joint system
cannot be evaluated without new data collection.

\emph{No raw video in DAIC-WOZ.} The publicly distributed DAIC-WOZ
release provides audio and transcripts; it does not provide the raw
facial video that the facial component would require. Facial modeling is
therefore grounded in FER2013 and AffectNet rather than in interview
data.

\emph{Linear projection.} The map $\phi$ is linear and may fail to
capture nonlinear cross-modal relationships.

\emph{Heuristic affect weights.} $w_V$ and $w_A$ are hand-set priors,
not estimated coefficients.

\emph{Distress ambiguity.} $D = (1-V)|A|$ cannot separate calm from
blunted affect; both yield $D \approx 0$.

\emph{Population distribution.} The baseline would be estimated from a
reference cohort, which makes $\mu$ and $\sigma$ unstable when that
cohort is small. The $|Z| > 2$ threshold assumes an approximately
Gaussian distribution and has not been clinically calibrated.

\emph{Cross-cultural validity.} The affect priors derive from a
predominantly Western annotation protocol and may not generalise.

\emph{Conflicting evidence handling.} Prediction-level divergence
detection is specified but not developed; conflict is captured only
partially, in feature space.

\emph{Environmental sensitivity.} Background noise and lighting would
affect speech and facial extraction.

\emph{Construct validity.} $C$ measures representational disagreement.
Whether this corresponds to the clinical construct of masking is an
open question that the proposed protocol can only partially address.

\emph{Not a diagnostic device.} Any deployed version would be a
screening aid and would not replace professional clinical assessment.

% =====================================================================
\section{Future Work}
\label{sec:future}

\begin{enumerate}
  \item \textbf{Execute the evaluation protocol.} Run the study
        specified in Section~\ref{sec:evaluation}, beginning with
        Hypotheses H1--H3 on component-level data and proceeding to H4
        once a joint multimodal cohort exists.
  \item \textbf{Prediction-level divergence detection.} Implement
        $D_{ET}$, $D_{EF}$, and $D_{TF}$ as formalised in
        Section~\ref{sec:divergence}, localising conflict to specific
        modality pairs.
  \item \textbf{Nonlinear projection.} Replace the linear $\phi$ with a
        shallow neural network and compare under identical splits.
  \item \textbf{Learned affect weights.} Refit $w_V$ and $w_A$ by
        regression against AffectNet dimensional annotations and
        evaluate with the concordance correlation coefficient.
  \item \textbf{Expressivity dynamics.} Validate session-level
        $\sigma_V$ and $\sigma_A$ as discriminators of flat versus calm
        affect.
  \item \textbf{Personalised baselines.} Estimate the blend coefficient
        $\kappa_i$ from longitudinal data and compare
        between-subject against within-subject normalisation.
  \item \textbf{Clinical calibration.} Calibrate the z-score threshold
        and the gate threshold $\tau_C$ against PHQ-9~\cite{kroenke2001phq9}
        and GAD-7~\cite{spitzer2006gad7} outcomes using ROC analysis.
  \item \textbf{Uncertainty and attribution.} Add Bayesian uncertainty
        estimation and SHAP-value attribution.
  \item \textbf{Privacy-preserving deployment.} Explore federated
        learning and on-device inference with clinical partners.
\end{enumerate}

% =====================================================================
\section{Conclusion}

This paper presented a conceptual framework for early mental health risk
detection that treats explicit--implicit inconsistency and
population-relative calibration as first-class concerns. The framework
combines an eight-module FSM interview engine grouped into five thematic
domains, a dual elicitation consistency model that quantifies
disagreement between self-report and inferred behaviour, a dimensional
facial component grounded in Russell's
circumplex~\cite{russell1980circumplex} and
AffectNet~\cite{mollahosseini2017affectnet}, and a z-score baseline
layer that situates each session against a population distribution.

The principal design argument is that the consistency score should be
treated as a reliability signal rather than as an additive severity
term. Discrepancy has no inherent clinical direction, an additive term
double-counts the score when it also appears in the feature vector, and
adding it where the estimate is least trustworthy inflates exactly the
cases that should be treated with most caution. The framework therefore
specifies a gated decision rule,
\begin{equation}
\text{escalate}(i) =
\mathbb{1}\bigl[\,C_i > \tau_{C} \;\lor\; |Z_{R,i}| > 2
\;\lor\; |Z_{L,i}| > 2 \;\lor\; \text{safety}_i = 1\,\bigr],
\end{equation}
with the reported output being the pair $\bigl(\hat{r}_i,
\hat{\rho}(C_i)\bigr)$ and the escalation flag.

No experimental results are reported. The contribution is a
specification and a falsifiable evaluation plan, not evidence of
superior performance. Whether the framework works---in particular,
whether $C$ predicts error as the design assumes---remains an open
empirical question, and Section~\ref{sec:evaluation} states the
experiment that would answer it.

% =====================================================================
\bibliographystyle{ieeetr}
\bibliography{references}

\end{document}
